% SHARD-ID: CA-17-LATTICE-SYMMETRY-EXAMPLES-QUEUE
% SHARD-TITLE: Small Examples and the Lattice-Symmetry Research Queue
% SHARD-SUMMARY: Works through a sequence of small examples showing what the lattice-symmetry compiler can and cannot infer.
% SHARD-SUMMARY: Records boundary cases: black-box Hamiltonians, open chains, nonuniform spacing, periodic chains, square lattices, and Koo-Saleur modes.
% SHARD-SUMMARY: Lists the next source acquisitions and derivations needed before promoting the programme beyond proposal status.
% SHARD-KEYWORDS: examples, research queue, open chain, periodic chain, square lattice, Koo-Saleur, Weinberg, source acquisition

\section{Small Examples and the Lattice-Symmetry Research Queue}
\label{sec:lattice-symmetry-examples-queue}

\subsection{Status, Sources, and Dependencies}

\textbf{Status: worked examples and queue.}  The examples are deliberately
small.  They illustrate compiler behaviour, not new convergence theorems.

Dependencies: CA-10--CA-16 and CONVENTIONS.md (h), (j).

% Source: references/cft/Schottenloher2008/Schottenloher2008.md:4046 -- time evolution from a one-parameter unitary group.
% Source: references/cft/Schottenloher2008/Schottenloher2008.md:4117 -- P_0=H and P_j are momenta.
% Source: references/text/CFTFromLatticeFermions.txt:2182--2195 -- lattice Fourier modes of Hamiltonian density.
% Source: references/text/CFTFromLatticeFermions.txt:2233--2235 -- smeared KS approximants converge in the cited setting.
% Convention: CONVENTIONS.md (j) -- finite periodic Gaussian grids are symbol/eigenvalue checks until a periodic coordinate convention is fixed.
% Checked: test/runtests.jl, "Poincare vector-field brackets" and "nearest-neighbour boost current coefficients".

\subsection{Example 1: Black-Box Quantum System}

Input: a Hilbert space and a self-adjoint Hamiltonian \(H\).

Output: time translation \(U(t)=e^{-itH}\).

Failed request: infer boosts, rotations, or spatial translations.  There is no
lattice geometry, no position map, and no Hamiltonian density.  The compiler
must stop.

\subsection{Example 2: Open Uniform Chain}

Input:
\[
  H=\sum_{j=1}^N h_j,\qquad x_j=j,\qquad [h_j,h_k]=0\text{ for }|j-k|>1.
\]
Candidate boost:
\[
  K=\sum_j jh_j.
\]
Candidate momentum:
\[
  P=i[H,K]=\sum_{j=1}^{N-1}i[h_j,h_{j+1}].
\]
Status: checked finite derivation; no continuum convergence yet.

\subsection{Example 3: Nonuniform Chain}

Input: positions \(x_1<\cdots<x_N\).  The same derivation gives
\[
  P_x=i[H,K_x]
  =
  \sum_{j=1}^{N-1}(x_{j+1}-x_j)i[h_j,h_{j+1}].
\]
The coefficients are discrete spacings.  The Julia test pins this invariant for
uniform and nonuniform examples.

\subsection{Example 4: Periodic Chain}

Input: a circle with periodic density terms.

Problem: \(x_j=j\) is not a single-valued function on the circle.  A periodic
boost-like first moment requires a branch cut, a sawtooth, Fourier modes, or a
different construction.  The same warning applies to finite periodic
momentum-coordinate tests: a finite grid supports periodic symbols and Fourier
eigenvalue checks, not an exact differential operator \(X=i\partial_k\) with a
canonical commutator.  Koo--Saleur Fourier modes are the better sourced
periodic prototype.

\subsection{Example 5: Square Lattice}

Input: positions \(r=(m,n)\).  Scalar ramps \(f_x(r)=m\), \(f_y(r)=n\) give two
candidate momentum components via
\[
  P_a^{\mathrm{lat}}=i[H,\sum_r r^a h_r].
\]
Rotation about the origin should then be built as a position-dependent
translation from the two momentum-density components.  Status: proposal until an
orientation and density-localisation convention is fixed.

\subsection{Example 6: Koo--Saleur Modes}

Input: periodic one-dimensional critical Hamiltonian density.  Fourier modes of
the density and commutator corrections give candidates for \(L_k,\bar L_k\).
In the cited free-fermion/OAR setting, smeared approximants converge to smeared
Virasoro generators.  Status outside that setting: source-local precedent, not a
general theorem.

\subsection{Research Queue}

\[
\begin{array}{ll}
1.&\text{Register Weinberg Vol. 1 or an equivalent QFT source with OCR anchors.}\\
2.&\text{Promote/register the original Koo--Saleur source before citing it.}\\
3.&\text{Source or derive a general local energy-current theorem.}\\
4.&\text{Fix periodic branch/sawtooth/Fourier first-moment conventions.}\\
5.&\text{Define higher-dimensional oriented density supports.}\\
6.&\text{Add numerical examples checking commutator closure in small chains.}\\
7.&\text{Connect anyon-chain Hamiltonian densities to the same candidate layer.}
\end{array}
\]

\subsection{Landing Point}

The next compiler upgrade is therefore not only a Hilbert-space compiler and not
yet only a tensor-category compiler.  It is a lattice-QM compiler layer:
\[
  (\text{carrier},\mathcal A_N,\Lambda_N,\{h_r\},x,S_N)
  \longmapsto
  (\text{candidate bulk symmetries},\text{evidence obligations}).
\]
Only after those obligations are discharged can the continuum symmetry
representation be claimed.
